% ============================================================
% abstract.tex
% ~190 words as written; the \todo results bring it to ~200.
% No citations in the abstract (AAAI convention).
% ============================================================

\begin{abstract}
Comparing machine-labelled sentiment across social media platforms is
now routine: differences in classifier scores between platforms are
read as differences in public opinion. These comparisons rest on an
assumption that the classifier measures the same construct, on the
same scale, on every platform. To our knowledge, no study has tested
that assumption for machine-labelled constructs by asking whether the
platform is part of the measurement. We adapt the psychometric equivalence toolkit (generalizability theory, measurement invariance and differential item functioning, and score linking) to machine-labelled policy sentiment, treating the platform as a survey
mode. We apply it to \todo{approximately 279{,}000 posts; confirm totals
and the collection window after recount: per-platform sums reach
roughly 395,000 and the results figures classify 255,446 items} about
UK economic policy from Twitter/X, Reddit, YouTube, and Quora
\todo{collection window, pending recount}, labelled by a fine-tuned BERT classifier and validated
against a probability-sampled human-annotated criterion \todo{to be
drawn; the dissertation-stage criterion was 200 balanced posts, kappa
approximately 0.78}. A same-event design controls topic and time,
pre-identified fiscal events centred on the September 2022 mini-budget
provide the linking anchors, and replication with a non-lexicon
transformer separates platform mode effects from lexicon artefacts.
\todo{One-sentence headline invariance result.} \todo{One-sentence
naive pooled index versus measurement-adjusted index result.} The protocol lets a researcher test whether pooling platforms is defensible in a given corpus, or whether platform scores must be linked before any average is taken.
\end{abstract}
